%% Field-projection mix-net proofs for BGV -- draft in ACM CCS (acmart sigconf) format.
%% Build:  latexmk -pdf main.tex
\documentclass[sigconf,anonymous]{acmart}

%% Submission settings: no copyright block, no ACM reference format.
\setcopyright{none}
\settopmatter{printacmref=false}
\renewcommand\footnotetextcopyrightpermission[1]{}
\acmConference[Draft]{Draft}{November 2026}{}
\acmYear{2026}
\acmDOI{}
\acmISBN{}

\usepackage{enumitem}
\usepackage{multirow}

%% ---------------------------------------------------------------------------------------
%% Notation
%% ---------------------------------------------------------------------------------------
\newcommand{\Z}{\mathbb{Z}}
\newcommand{\Zq}{\mathbb{Z}_q}
\newcommand{\Rq}{R_q}
\newcommand{\F}{\mathbb{F}}
\newcommand{\cube}[1]{\{0,1\}^{#1}}
\newcommand{\eq}{\mathsf{eq}}
\newcommand{\Enc}{\mathsf{Enc}}
\newcommand{\EncZ}{\mathsf{Enc}_0}
\newcommand{\coeff}{\mathsf{coeff}}
\newcommand{\Mpk}{M_{\mathsf{pk}}}
\newcommand{\Rmix}{\mathcal{R}_{\mathsf{mix}}}
\newcommand{\Rdec}{\mathcal{R}_{\mathsf{dec}}}
\newcommand{\RS}{\mathsf{RS}}
\newcommand{\ip}[2]{\langle #1, #2\rangle}
\newcommand{\mle}[1]{\widetilde{#1}}
\newcommand{\msk}[1]{#1^{\star}}
\newcommand{\epsprox}{\varepsilon_{\mathsf{prox}}}
\newcommand{\epsrbr}{\varepsilon_{\mathsf{rbr}}}
\newcommand{\nq}{n_{\mathsf{q}}}
\newcommand{\uh}{\hat{u}}
\newcommand{\hh}{\hat{h}}
\newcommand{\zh}{\hat{z}}
\newcommand{\ch}{\hat{c}}


%% Measurements of the extracted shuffle prover and verifier (scratchpad xbench, M5 Max, one core).
\newcommand{\xbProveA}{0.30}  % 256 ballots: prover seconds per ballot (77.4 s in total)
\newcommand{\xbVerA}{5.6}    % 256 ballots: verifier seconds in total
\newcommand{\xbSizeA}{2.19}  % 256 ballots: proof MB (2194984 bytes)
\begin{document}

\title{Leaving the Ring: Hash-Based Proofs of Shuffle and Decryption for BGV Mix-Nets}

\author{Anonymous Author(s)}
\affiliation{\institution{Anonymous Institution}\country{}}

\begin{abstract}
A verifiable mix-net over lattice ciphertexts needs two zero-knowledge proofs. A mix server must
prove that it permuted and re-randomised its inputs, and a decryption server must prove that it
applied its key share correctly. The most efficient published construction, by Aranha, Baum,
Gj{\o}steen and Silde (CCS 2023), proves the shuffle with a polynomial identity over the ciphertext
ring $\Rq=\Zq[X]/(X^N+1)$. Bootle, Lyubashevsky and Merino-Gallardo (PKC 2025) showed that this
argument is unsound: $\Rq$ has zero divisors and no unique factorisation.

We give proofs for the same BGV mix-net relation, with the same keys and ciphertexts, that never
test an identity over the ring. The relation is $\Zq$-linear in the coefficient vectors. The
verifier therefore \emph{projects} every ciphertext to a single element of the field
$\F=\mathrm{GF}(q^2)$ with a random linear form, and the shuffle becomes a statement about two
lists of field elements. The approach is modular: after the projection, what remains is a linear
map, a permutation and a range check over $\F$, and any commit-and-prove system for these
relations can be used. We instantiate it with a tagged logarithmic-derivative argument, three
masked sumchecks, and a zero-knowledge opening of a witness that is encoded with a Reed--Solomon
code and committed with Merkle trees. The proofs are non-interactive and statistically
zero-knowledge. They are knowledge-sound in the random-oracle model, and the proof system uses no
lattice assumption. The Fiat--Shamir analysis goes through round-by-round soundness, which also
covers quantum adversaries. The decryption proof contains an exact range proof for the flooding
noise.

We machine-checked the following in the Rocq prover, in about $26{,}500$ lines and with no axioms:
the soundness theorems and honest-verifier zero knowledge of both proofs, and the correctness of
executable verifiers and provers. For both proofs, zero knowledge holds for the executable provers
themselves: the field elements of every proof they output form the honest transcript of the model
for which we prove the simulator. The cited proximity lemma for Reed--Solomon codes is a
hypothesis, and the Fiat--Shamir step is argued on paper. The shuffle prover and verifier,
extracted to OCaml, run at full parameters. For ring degree $4096$, a $78$-bit modulus and
$128$-bit security, they prove a shuffle of $256$ ballots in \xbProveA\,s per ballot on one core
and verify it in \xbVerA\,s. The proof has \xbSizeA\,MB.
\end{abstract}

\begin{CCSXML}
<ccs2012>
 <concept>
  <concept_id>10002978.10002979.10002981</concept_id>
  <concept_desc>Security and privacy~Public key (asymmetric) techniques</concept_desc>
  <concept_significance>500</concept_significance>
 </concept>
 <concept>
  <concept_id>10002978.10002991.10002995</concept_id>
  <concept_desc>Security and privacy~Privacy-preserving protocols</concept_desc>
  <concept_significance>500</concept_significance>
 </concept>
 <concept>
  <concept_id>10002978.10002979.10002980</concept_id>
  <concept_desc>Security and privacy~Key management</concept_desc>
  <concept_significance>100</concept_significance>
 </concept>
</ccs2012>
\end{CCSXML}
\ccsdesc[500]{Security and privacy~Public key (asymmetric) techniques}
\ccsdesc[500]{Security and privacy~Privacy-preserving protocols}

\keywords{mix-nets, proof of shuffle, verifiable decryption, BGV, post-quantum, electronic voting,
sumcheck, interactive oracle proofs}

\maketitle

%% =======================================================================================
\section{Introduction}
%% =======================================================================================

A mix-net~\cite{Chaum81} takes a list of ciphertexts and outputs a list that decrypts to the same
plaintexts in an order that cannot be linked to the input. In electronic voting this separates
ballots from voters. Each mix server permutes and re-randomises the list, and a set of decryption
servers then opens the result. Every server must prove that it behaved correctly without revealing
its permutation or its key share. For discrete-logarithm ciphertexts, such proofs have been known and
deployed for two decades~\cite{Neff01,BG12}. Lattice ciphertexts are the candidates for ballots
that must stay secret against quantum adversaries, and for them these proofs are still a research
topic~\cite{CMM19,Strand18,ABGSctrsa21,FWK21,HMS21,ABGS23,BLM25}.

The most efficient lattice mix-net with a published implementation is that of Aranha, Baum,
Gj{\o}steen and Silde~\cite{ABGS23}. It uses BGV encryption~\cite{BGV12} over
$\Rq=\Zq[X]/(X^N+1)$ together with lattice-based commitments and proofs. Bootle, Lyubashevsky and
Merino-Gallardo~\cite{BLM25} later found that its shuffle argument, which goes back
to~\cite{ABGSctrsa21}, is unsound. The argument tests a polynomial identity with coefficients in
$\Rq$ and then reasons as if $\Rq$ had unique factorisation. For the moduli in use, $\Rq$ splits
into a product of smaller rings, its CRT \emph{slots}. It therefore has zero divisors, and a
polynomial over it can vanish at a random point without being zero. A cheating server can swap a
single slot between two ciphertexts. The result decrypts to different votes and still satisfies
the tested identity.

\paragraph{Our approach.}
We keep the encryption scheme, the keys and the ciphertexts of~\cite{ABGS23} and replace the two
proofs. Our starting point is that the mix-net relation does not need the ring structure. Write
$\EncZ(w)=(ar+pe_1,\;br+pe_2)$ for an encryption of zero under the public key $(a,b)$ with
randomness $w=(r,e_1,e_2)$. Output $j$ of a correct shuffle is
\[
  \ch_j \;=\; c_{\sigma(j)} + \EncZ(w_{\sigma(j)}),
\]
and multiplication by the \emph{public} elements $a$ and $b$ is a linear map on coefficient vectors.
The relation is therefore a system of linear equations over $\Zq$ in the entries of the $w_i$, plus
two non-linear side conditions: $\sigma$ is a permutation, and the $w_i$ are short. Linear equations
over $\Zq$ can be checked over any field that contains $\Zq$. We use $\F=\mathrm{GF}(q^2)$, which
has about $2^{156}$ elements for the parameters of~\cite{ABGS23}, and we never test an identity over
$\Rq$. The slot attack has no counterpart, because the argument never looks at slots.

The proof combines four techniques.
\begin{description}[leftmargin=1em,itemsep=2pt,topsep=3pt]
\item[Field projection.] The verifier draws $\rho\in\F^{2N}$ and maps each ciphertext $c$ to
  $\ip{\rho}{\coeff(c)}\in\F$. The projection of a re-randomised input $c_i+\EncZ(w_i)$ is a public
  value plus the inner product of the secret $w_i$ with a public vector. A shuffle of ciphertexts
  becomes a permutation of $\tau$ field elements, and two different ciphertexts have the same
  projection with probability $1/|\F|$. A related compression of ring elements to field elements, by
  evaluation at a random point, is used by Fiore, Nitulescu and Pointcheval~\cite{FNP20}.
\item[Tagged logarithmic derivative.] We tag input $i$ with its index $i$ and output $j$ with the
  committed value $\sigma_j$. The prover shows
  $\sum_i 1/(z_i+\beta i-\alpha)=\sum_j 1/(\zh_j+\beta\sigma_j-\alpha)$ for random
  $\alpha,\beta$~\cite{Hab22}. Equality of the tagged multisets forces $\sigma$ to be a
  permutation. Unlike a grand-product argument, this reveals no value that depends on $\sigma$.
\item[Masked sumchecks.] Three sumchecks~\cite{LFKN92} reduce the inverse relations, the linear map
  from the witness to the projections, and the shortness condition $w^3=w$ to a few evaluations of
  committed tables at random points. Each sumcheck is made zero-knowledge with the masking technique
  of~\cite{CFS17,Libra}.
\item[Hash-based commitment.] The prover commits to the witness row by row, with a Reed--Solomon code
  and a Merkle tree over the columns, as in Ligero~\cite{AHIV17}. We give a zero-knowledge opening
  for batches of linear claims, with two parameter regimes for the proximity
  test~\cite{DP23,BCIKS20}.
\end{description}
The proof of correct partial decryption follows the same pattern. Its relation $t_i=su_i+pE_i$ is
linear in the key share $s$ and the noise $E_i$. The bound on $E_i$ is proven exactly, with a digit
decomposition and one more sumcheck.

\paragraph{Contributions.}
\begin{itemize}[leftmargin=1.2em,itemsep=2pt,topsep=3pt]
\item A proof of shuffle and a proof of correct partial decryption for the BGV mix-net relation
  of~\cite{ABGS23}, with all algebra over a field (Sections~\ref{sec:shuffle} and~\ref{sec:dec}).
  They replace the existing proofs without changing the encryption scheme or the ciphertexts.
\item Security proofs: knowledge soundness with an explicit error budget (Theorem~\ref{thm:sound});
  honest-verifier zero knowledge with a simulator that is exact except for unopened hash values
  (Theorem~\ref{thm:zk}); and a Fiat--Shamir analysis through round-by-round
  soundness~\cite{BCS16,Hol19,CMS19} (Section~\ref{sec:fs}). Soundness is statistical in the
  random-oracle model, and no lattice assumption enters the proof system.
\item A mechanisation in Rocq~\cite{Rocq} (Section~\ref{sec:mech}) of knowledge soundness and
  honest-verifier zero knowledge of both protocols, for the exact row layout of the
  implementation, with the cited proximity lemma as the only hypothesis. It found a few unstated
  conditions and one imprecision, and confirmed every error bound (Section~\ref{sec:findings}).
\item Executable verifiers and provers for both proofs, written in Rocq and extracted to
  OCaml~\cite{Letouzey02}. We prove the extracted provers complete and faithful to the
  zero-knowledge model, and the algebraic part of the extracted verifiers sound with respect to
  the interactive model (Section~\ref{sec:mech}).
\end{itemize}

\paragraph{Status of the claims.}
The following results are machine-checked: Theorems~\ref{thm:sound} and~\ref{thm:dec} (soundness),
Lemmas~\ref{lem:logder}, \ref{lem:cons} and~\ref{lem:hide}, Theorem~\ref{thm:zk}, the
zero-knowledge part of Theorem~\ref{thm:dec}, and the extension of the line lemma in
Appendix~\ref{app:prox}. The line lemma and the proximity gap of~\cite{BCIKS20} are cited, not
mechanised. The machine-checked statements are about the interactive protocols, with the committed
matrices as oracles. Merkle trees, the binding of the roots and the Fiat--Shamir transformation are
outside the model. We prove the round-by-round analysis behind the non-interactive statement
(Section~\ref{sec:fs}) on paper, with a state function written out message by message; its
per-message transitions are the mechanised counts. The code is a research prototype.
Section~\ref{sec:limits} lists the open items.

%% =======================================================================================
\section{Preliminaries}
\label{sec:prelim}
%% =======================================================================================

\subsection{BGV and the mix-net relations}

Let $N$ be a power of two, $q$ a prime with $q\equiv 1 \pmod{2N}$, $\Rq=\Zq[X]/(X^N+1)$, and $p$ a
small plaintext modulus. For $x\in\Rq$ we write $\coeff(x)\in\Zq^N$ for its coefficient vector,
and we extend this notation to tuples by concatenation. An element is \emph{ternary} if all its
coefficients lie in $\{-1,0,1\}$. Following the parameter choice of~\cite{ABGS23}, our
implementation uses $N=4096$, $p=2$ and $q=2^{78}+100663297$.

A BGV public key is a pair $(a,b)\in\Rq^2$ with $b=as+pe$ for ternary $s,e$. A ciphertext is a
pair $c=(u,v)$, and
\[
  \Enc(m;w)=(ar+pe_1,\; br+pe_2+m),\qquad w=(r,e_1,e_2)\ \text{ternary}.
\]
We write $\EncZ(w)=\Enc(0;w)$. Decryption computes $(v-su \bmod q)\bmod p$. In distributed
decryption, each decryption server holds a share $s$ with public $b=as+pe$ and publishes
$t_i=su_i+pE_i$ for every ciphertext, where $E_i$ is flooding noise.

\begin{definition}[Shuffle relation]
A statement consists of a public key and two lists of $\tau$ ciphertexts,
$x=(a,b,(c_i)_{i<\tau},(\ch_j)_{j<\tau})$. A witness is a permutation $\sigma$ of
$\{0,\dots,\tau-1\}$ together with ternary elements $w_0,\dots,w_{\tau-1}$ of $\Rq^3$ such that
$\ch_j=c_{\sigma(j)}+\EncZ(w_{\sigma(j)})$ for all $j$. We write $\Rmix$ for this relation.
\end{definition}

\begin{definition}[Decryption relation]
A statement consists of $a,b\in\Rq$ and of elements $u_i,t_i\in\Rq$ for $i<\tau$. A witness is a
pair of ternary $s,e\in\Rq$ and elements $E_i\in\Rq$ with integer coefficients in $[-B/2,B/2)$
such that
\[
  \text{(K)}\;\; b=as+pe \qquad\text{and}\qquad \text{(D)}\;\; t_i=su_i+pE_i \;\;\text{for all } i .
\]
We write $\Rdec$ for this relation.
\end{definition}

\paragraph{The linear form.}
Since $a$ and $b$ are public, there is a public matrix $\Mpk\in\Zq^{2N\times 3N}$ with
\begin{equation}\label{eq:linear}
  \coeff(\EncZ(w)) \;=\; \Mpk\,\coeff(w)\qquad\text{for all } w\in\Rq^3 .
\end{equation}
Its blocks are the negacyclic multiplication matrices of $a$ and $b$ and multiples of the identity.
Equation~\eqref{eq:linear} is the only property of $\Rq$ that our proofs use.

\subsection{Field, tables and multilinear extensions}

Let $\F=\Zq[t]/(t^2-3)$. This is a field with $q^2\approx 2^{156}$ elements, because $3$ is not a
square modulo our $q$. All verifier challenges are elements of $\F$. Committed data are elements of
$\Zq\subset\F$, and a committed $\F$-element is stored as two $\Zq$-elements.

A \emph{table} $T$ of length $2^m$ is a function on $\cube{m}$; variable $0$ is the least
significant bit of the index. Its multilinear extension is
$\mle{T}(x)=\sum_{b\in\cube{m}}T(b)\,\eq(b,x)$, where
$\eq(t,x)=\prod_i\bigl(t_ix_i+(1-t_i)(1-x_i)\bigr)$. Let $Z(x)=\prod_i x_i(1-x_i)$, which vanishes on
$\cube{m}$. For a uniformly random \emph{mask coefficient} $a_T\in\F$, the \emph{masked extension}
\[
  \msk{T}(x)=\mle{T}(x)+a_T\,Z(x)
\]
agrees with $T$ on the cube, and $\msk{T}(r)$ is uniform in $\F$ at every point $r$ with
$Z(r)\ne 0$.

We write the witness of $\Rmix$ as three tables $W_r,W_{e_1},W_{e_2}$ of length $\tau N$, with
coefficient $y$ of ballot $i$ at index $y+Ni$, and a table $\sigma$ of length $\tau$. We assume that
$\tau$ is a power of two and set $\ell=\log\tau$ and $m=\ell+\log N$.

\subsection{Sumcheck}

The sumcheck protocol~\cite{LFKN92} reduces a claim $\sum_{x\in\cube{m}}g(x)=S$ to one evaluation of
$g$ at a random point. In round $j$ the prover fixes the first $j-1$ variables to the challenges,
sums over the remaining ones, and sends the resulting univariate polynomial. The verifier checks
that its values at $0$ and $1$ add up to the running claim and replies with a random $r_j\in\F$.
If $g$ has degree $d_j$ in variable $j$, a false claim survives with probability at most
$\sum_j d_j/|\F|$. Our prover sends the values at $0,2,\dots,d_j$; the value at $1$ follows from the
running claim.

To hide the round messages, the prover adds a random polynomial to the
summand~\cite{CFS17,Libra}. The mask is a sum of univariate polynomials,
\[
  g(x)=a_0+\sum_{j}\sum_{t=1}^{d_j}c_{j,t}\,x_j^{\,t},
\]
and the prover commits to its coefficients together with its sum $G$ over the cube. The verifier
sends a random $\rho\in\F$, and the parties run the sumcheck on $f+\rho g$ with claim $S+\rho G$.

\subsection{Proof models}

We first analyse the protocols as interactive oracle proofs~\cite{BCS16}: the long messages of the
prover are oracles, and the verifier is public-coin. The BCS transformation~\cite{BCS16} replaces
oracles by Merkle roots and challenges by hashes of the transcript. Its soundness in the
random-oracle model is governed by round-by-round soundness~\cite{Hol19}, and the same quantity
bounds soundness against quantum adversaries~\cite{CMS19}.

%% =======================================================================================
\section{Technical overview}
\label{sec:overview}
%% =======================================================================================

\paragraph{Why not the ring.}
A proof of shuffle in the style of Neff~\cite{Neff01} shows that two lists are equal as multisets by
comparing $\prod_i(x_i-\alpha)$ and $\prod_j(\hat{x}_j-\alpha)$ at a random $\alpha$. Over a field,
two different polynomials of degree $\tau$ agree on at most $\tau$ points. Over $\Rq$ this fails.
For $q\equiv 1\pmod{2N}$ the ring is isomorphic to $\Zq^N$, and the comparison happens slot by slot,
so the two products can agree in every slot although the lists differ. Swapping one slot between
two output ciphertexts changes both plaintexts but leaves every slot-wise multiset unchanged. Our
extracted verifier rejects statements modified in this way (Section~\ref{sec:eval}).

\paragraph{Projection.}
We want an identity test over a field. By~\eqref{eq:linear}, the re-randomised input
$c_i+\EncZ(w_i)$ has coefficient vector $\coeff(c_i)+\Mpk\coeff(w_i)$. For a random
$\rho=(\rho_u,\rho_v)\in\F^{2N}$ define
\begin{align*}
  z_i &= \ip{\rho}{\coeff(c_i)}+\ip{\rho'}{\coeff(w_i)}, & \rho' &=\rho\,\Mpk\in\F^{3N},\\
  \zh_j &= \ip{\rho}{\coeff(\ch_j)} .
\end{align*}
The verifier computes $\rho'$, the public parts of the $z_i$, and all $\zh_j$ by itself. Two
distinct vectors in $\Zq^{2N}$ have the same projection with probability $1/|\F|$. With
overwhelming probability, the shuffle relation therefore holds exactly when $\zh_j=z_{\sigma(j)}$
for all $j$. The projection mixes all coefficients, and hence all slots, into one field element.

\paragraph{Permutation without leakage.}
To show that $(\zh_j)$ is $(z_i)$ reordered by a \emph{permutation}, we attach tags. Let
\[
  h_i=z_i+\beta i-\alpha,\qquad \hh_j=\zh_j+\beta\sigma_j-\alpha ,
\]
for random $\alpha,\beta\in\F$, where $\sigma_j$ is the committed table. If
$\sum_i 1/h_i=\sum_j 1/\hh_j$, then with high probability the multiset of pairs $(z_i,i)$ equals
the multiset of pairs $(\zh_j,\sigma_j)$ (Lemma~\ref{lem:logder}). The tags on the left are
$0,\dots,\tau-1$, each exactly once, so $\sigma$ is a bijection and $\zh_j=z_{\sigma(j)}$. The
prover commits to the inverses $u=1/h$ and $\uh=1/\hh$ and proves $u\cdot h=1$, $\uh\cdot\hh=1$ and
$\sum u=\sum\uh$. An earlier version of our design compared products and sent their common value in
the clear. That value is a function of the secret permutation, and with two ballots it reveals the
permutation. In the logarithmic-derivative form, the prover sends only masked evaluations.

\paragraph{Three sumchecks and one opening.}
Sumcheck~A proves the three facts about $u$ and $\uh$. It ends with claimed evaluations of
$\msk{u},\msk{\uh},\msk{\sigma}$ and $\msk{h}$ at a point $r\in\F^{\ell}$. The prover never commits to
the table $h$: it is a linear image of the committed $W$, and sumcheck~B reduces the claimed value
$\msk{h}(r)$ to one linear claim about $W$. Sumcheck~C proves $w^3-w=0$ for all $3\tau N$ entries of
$W$, which over $\Zq$ with $q>3$ means $w\in\{-1,0,1\}$. What remains are nine linear claims about
committed rows, and one batched opening settles them.

\paragraph{Cost.}
The sumcheck part of a proof takes a few kilobytes. The opening dominates the size: it sends a
constant number of row combinations and $\nq\approx 190$ to $310$ columns of each committed matrix.
Splitting a table of $3\tau N$ symbols into rows of $C$ symbols gives $3\tau N/C$ rows, so the size
is $O(C+\nq\,\tau N/C)$ field elements, with minimum $O(\sqrt{\nq\,\tau N})$. The proof size per
ballot therefore shrinks as the number of ballots grows.

\paragraph{A general template.}
The projection is the only step that uses the ring. After it, the prover must show a statement over
a field: committed vectors $W$ with entries in $\Zq$ and a committed table $\sigma$ such that, for
the verifier's $\rho$ drawn after the commitment, the values
$z_i=\ip{\rho}{\coeff(c_i)}+\ip{\rho'}{\coeff(w_i)}$ are the public $\zh_j$ reordered by $\sigma$, and
every entry of $W$ lies in $\{-1,0,1\}$. These are a linear map, a permutation and a range check, and
any commit-and-prove system for them can replace our choices in steps 2 to 4. The permutation can
also be shown by a grand-product argument~\cite{Neff01}, provided its value stays masked, since the
product itself reveals $\sigma$. The reduction to evaluations can use univariate quotient checks
instead of multilinear sumchecks, as long as the commitment matches. The commitment can be a
polynomial commitment with polylogarithmic openings, such as BaseFold~\cite{BaseFold}, which would
make the proof much smaller than our $O(\sqrt{\nq\,\tau N})$ opening, or a lattice-based
one~\cite{LaBRADOR,NS24}, at the price of computational soundness. Three constraints apply to every
choice. The witness must be committed before $\rho$, $\alpha$ and $\beta$ are drawn. The data live in
$\Zq$ while the challenges need a larger field such as $\mathrm{GF}(q^2)$, so a system tied to a
different field would have to emulate arithmetic modulo $q$. And keeping the proof system free of
lattice assumptions rules out pairing-based and discrete-logarithm components. Our components are
hash-based and work over $\Zq$ directly. The security argument is modular in the same way, but the
mechanisation of Section~\ref{sec:mech} covers this instantiation only.

\paragraph{Other encryption schemes.}
The proofs use three properties of BGV. Re-randomisation adds an encryption of zero that is a
public linear map of the randomness, which is~\eqref{eq:linear}. The randomness is short in a way
that a low-degree check per coefficient can express. And the modulus $q$ is prime, so that $\Zq$
lies in a field from which the challenges are drawn. The ring enters only through the matrix
$\Mpk$: it need not split, and it need not be cyclotomic. Other schemes with these properties fit
the same template. BFV~\cite{FV12} has the ciphertext structure of BGV. For encryption from plain
or module LWE~\cite{Regev05}, without ciphertext compression, the public key is a matrix and
$\Mpk$ is built from it directly. Randomness from a wider range needs a check of higher degree or
the digit decomposition of Section~\ref{sec:dec}. Some cases need more than a change of
parameters. A composite or power-of-two modulus gives no field. A small prime modulus calls for a
challenge field of high extension degree and for a different row code, since $\Zq$ has too few
evaluation points. Rounding, as in ciphertext compression, is not linear and would need a range
proof for the rounding error. We work out and mechanise the BGV case only.

%% =======================================================================================
\section{Commitment and zero-knowledge opening}
\label{sec:pcs}
%% =======================================================================================

\subsection{Row code and commitment}

A row holds $C$ data symbols of $\Zq$ followed by $\nq$ uniformly random \emph{pads}. The
$k=C+\nq$ symbols are the coefficients of a polynomial, and the codeword is the list of its values
at the $n$-th roots of unity in $\Zq$. This is a Reed--Solomon code with distance $d=n-k+1$. The
length is $n=C/\varrho$ for a nominal rate $\varrho\in\{1/4,1/2\}$, and it is doubled if the pads
would push the actual rate above $2\varrho$. Fix a divisor $\Delta\in\{3,2\}$, the \emph{proximity
regime}, and let
\[
  e=\lfloor (d-1)/\Delta\rfloor,\qquad
  \nq=\min\{\,t : (1-e/n)^{t}\le 2^{-\lambda}\,\},
\]
with $\lambda=128$. A batch of rows forms a matrix. Leaf $j$ of its Merkle tree is the SHA-256 hash
of a fresh $16$-byte salt and column $j$, and the commitment is the root. A matrix may be committed
in several batches, each with its own root.

\subsection{Linear claims and the opening protocol}

A \emph{claim} is an equation $\sum_{t}\mathit{coef}_t\cdot\ip{\mathit{row}_t}{L_t}=\mathit{value}$
that relates committed rows, public vectors $L_t\in\F^{C}$ and a value in $\F$. An evaluation of a
multilinear extension at a point has this form: if the table is split into rows, $L$ is the $\eq$
vector of the low variables, and the coefficients are the $\eq$ values of the high variables. Rows
that are paired with the same vector $L$ form a \emph{slot} $v$ with vector $L_v$. Slots come in two
kinds. A \emph{tensor} slot has one fixed vector $L_v$ and many rows, as for a table split into rows.
A \emph{functional} slot has one row and, for each claim $i$, its own vector $L_{v,i}$; after
batching, its vector is $L_v=\sum_i\kappa_i\,\mathit{coef}_{i}\,L_{v,i}$. Functional slots read
single entries or short combinations of one row, for instance a mask coefficient in the mask table
or the value of a sumcheck mask at a point. For every slot the prover has committed to a uniformly
random mask row $m_v\in\F^{k}$, and for every matrix to one more mask row $m_{\mathsf{prox}}$. A mask
row over $\F$ is stored as two rows over $\Zq$, one per coordinate.

\begin{enumerate}[leftmargin=1.5em,itemsep=1pt,topsep=3pt]
\item The verifier sends $\kappa_i\in\F$ for every claim $i$. This defines, for each slot,
  $u_v=\sum_i\kappa_i\sum_{t\in v}\mathit{coef}_{i,t}\,\mathit{row}_{i,t}$ and the batched value
  $V=\sum_i\kappa_i\,\mathit{value}_i$.
\item The prover sends $\mu=\sum_v\ip{m_v}{L_v}$.
\item The verifier sends $c\in\F$, and $\gamma_i\in\F$ for every row $i$ of every matrix.
\item The prover sends $\bar u_v=m_v+c\,u_v$ for every slot, with all $k$ coordinates, and
  $p=m_{\mathsf{prox}}+\sum_i\gamma_i\,\mathit{row}_i$ for every matrix.
\item The verifier checks $\sum_v\ip{\bar u_v}{L_v}=\mu+cV$, where the inner product uses the first
  $C$ coordinates.
\item The verifier sends $\nq$ distinct column indices per matrix, and the prover opens these
  columns with their salts and Merkle paths.
\item For every opened column $j$, the verifier encodes $\bar u_v$ and $p$ and checks that their
  $j$-th symbols equal the same combinations of the opened entries:
  \begin{align*}
   \RS(\bar u_v)[j]&=\mathrm{col}_j[m_v]+c\sum_i\kappa_i\sum_{t\in v}\mathit{coef}_{i,t}\,\mathrm{col}_j[\mathit{row}_{i,t}],\\
   \RS(p)[j]&=\mathrm{col}_j[m_{\mathsf{prox}}]+\sum_i\gamma_i\,\mathrm{col}_j[i],
  \end{align*}
  where $\RS(\cdot)$ is the encoding map and $\mathrm{col}_j[\cdot]$ is the opened entry of a row.
\end{enumerate}

\subsection{Properties}

\begin{lemma}[Binding]\label{lem:bind}
In the random-oracle model, a Merkle root fixes at most one value per leaf, except with probability
$t^2/2^{256}$ for $t$ oracle queries~\cite{BCS16}.
\end{lemma}

\begin{lemma}[Proximity]\label{lem:prox}
Let $M$ be the rows of a matrix other than $m_{\mathsf{prox}}$. Suppose there is no set of $n-e$
columns on which every row of $M$ agrees with a codeword. Then for every $m_{\mathsf{prox}}$, the
combination $p_\gamma=m_{\mathsf{prox}}+\sum_i\gamma_i\,\mathit{row}_i$ satisfies
\[
  \Pr_{\gamma}\bigl[\,p_\gamma \text{ is within distance } e \text{ of the code over } \F\,\bigr]\le\epsprox ,
\]
where $\epsprox=(e+2)/|\F|$ for $\Delta=3$ and $\epsprox=n/|\F|$ for $\Delta=2$.
\end{lemma}

For $\Delta=2$, this is the proximity gap for Reed--Solomon codes at the unique-decoding radius,
with correlated agreement over affine spaces~\cite[Thm.~1.6]{BCIKS20}. For $\Delta=3$, it needs
only an elementary statement about lines that holds for every linear code and is attributed to Roth
and Z\'emor~\cite{AHIV17}. We use it as stated in~\cite[Thm.~2.1]{DP23} and extend it from lines to
affine spaces in Appendix~\ref{app:prox}. Rows are committed over $\Zq$ and tested over $\F$. A row
over $\Zq$ that agrees with an $\F$-codeword on more than $k$ columns agrees there with a
$\Zq$-codeword, namely the $\Zq$-component of the $\F$-codeword. The choice of regime trades
assumptions for size: $\Delta=3$ rests on the elementary lemma, while $\Delta=2$ needs fewer queries
and relies on the deeper theorem.

Soundness does not need this remark about the base field. The mechanised proof of
Theorem~\ref{thm:sound} decodes the committed rows over $\F$ and extracts a witness over $\F$. The
shortness sumcheck then shows $w^3=w$, which forces every entry into $\{-1,0,1\}$, and the
multiset step forces $\sigma$ to be a permutation. No step of the argument uses that the decoded
entries lie in $\Zq$.

\begin{lemma}[Consistency]\label{lem:cons}
Suppose the rows of $M$ agree with codewords outside a set of at most $e$ columns, and some
$\bar u_v$ differs from $m_v+c\,u_v$ computed on the decoded rows. Then a uniformly random column
passes the check of step~7 with probability at most $1-(d-e)/n$.
\end{lemma}
\begin{proof}
The encoding of $\bar u_v$ and the decoded combination are distinct codewords, so they differ in at
least $d$ columns. At most $e$ of these columns are ones where the committed rows differ from their
decodings.
\end{proof}

\begin{lemma}[Hiding]\label{lem:hide}
For every set $J$ of at most $\nq$ columns, the entries of a row at $J$ are uniform and independent
of its data symbols. More precisely, for $|J|=\nq$ distinct columns, the map from the $\nq$ pads to
the entries at $J$ is a bijection, whatever the data symbols.
\end{lemma}
\begin{proof}
For pads $\pi_t$, the symbol at column $j$ is the data part plus
$\omega^{jC}\sum_{t<\nq}\pi_t\,\omega^{jt}$. Restricted to $J$, this is an invertible diagonal matrix
times a Vandermonde matrix with distinct nodes, applied to uniform pads.
\end{proof}

Lemmas~\ref{lem:cons} and~\ref{lem:hide} are machine-checked, and so is the soundness of the
opening as a whole, for two matrices and both kinds of slot. If a claimed value is not the value of
its claim on the decoded rows, then the value check and all column checks pass for at most the
fraction of challenges given by the last four lines of Table~\ref{tab:budget}. One detail is sharper
than a union bound over the matrices. For fixed earlier challenges, the queries need to catch at
most one matrix, so the query term is the larger of the two matrices' terms rather than their sum.

%% =======================================================================================
\section{The shuffle proof}
\label{sec:shuffle}
%% =======================================================================================

\subsection{Protocol}

The prover commits to two matrices. The \emph{big} matrix holds $W_r,W_{e_1},W_{e_2}$ in rows of
$C_1$ symbols, plus three mask rows over $\F$. The \emph{small} matrix has rows of $C_2$ symbols and
is committed in two batches: first $\sigma$ and three mask rows, then, after the first challenges,
$u$, $\uh$ and one row that holds the \emph{mask table}. The mask table contains the mask
coefficients $a_u,a_{\uh},a_\sigma,a_h,a_{w_1},a_{w_2},a_{w_3}$ and the coefficients of the three
sumcheck masks $g_A,g_B,g_C$. Figure~\ref{fig:mix} gives the protocol.

\begin{figure*}[t]
\small
\fbox{\begin{minipage}{0.975\textwidth}
\textbf{Statement} $x=(a,b,(c_i),(\ch_j))$.\quad \textbf{Witness} $\sigma$, $W=(W_r,W_{e_1},W_{e_2})$.
\begin{enumerate}[leftmargin=1.6em,itemsep=2pt,topsep=3pt]
\item[\textbf{1.}] $P$ commits $W$ (big matrix) and $\sigma$ (small matrix, batch 1), each with its mask
  rows, and sends the two roots.
\item[\textbf{2.}] $V$ sends $\rho=(\rho_u,\rho_v)\in\F^{2N}$ and $\alpha,\beta\in\F$. Both compute
  $\rho'=\rho\,\Mpk$, $z^{\mathsf{pub}}_i=\ip{\rho}{\coeff(c_i)}$ and $\zh_j=\ip{\rho}{\coeff(\ch_j)}$.
  $P$ computes
  \[ z_i=z^{\mathsf{pub}}_i+\ip{\rho'}{\coeff(w_i)},\qquad h_i=z_i+\beta i-\alpha,\qquad
     \hh_j=\zh_j+\beta\sigma_j-\alpha,\qquad u_i=1/h_i,\qquad \uh_j=1/\hh_j . \]
\item[\textbf{3.}] $P$ commits $u$, $\uh$ and the mask table (small matrix, batch 2) and sends the
  root and the sums $G_A,G_B,G_C$ of $g_A,g_B,g_C$ over their cubes.
\item[\textbf{4.}] \emph{Sumcheck A} (permutation; $\ell$ variables, degree $3$, last round $5$).
  $V$ sends $t_c\in\F^{\ell}$, $k_1,k_2,k_3,\rho_A\in\F$. The parties run a sumcheck for
  \[ \sum_{x\in\cube{\ell}}\Bigl[\eq(t_c,x)\bigl(k_1\,\msk{u}\msk{h}+k_2\,\msk{\uh}\msk{\hh}\bigr)
       +k_3\bigl(\msk{u}-\msk{\uh}\bigr)+\rho_A\,g_A\Bigr](x)\;=\;k_1+k_2+\rho_A G_A ,
     \qquad \msk{\hh}(x)=\mle{\zh}(x)+\beta\,\msk{\sigma}(x)-\alpha . \]
  It ends at a point $r$. $P$ sends
  $(U,\hat U,s,H)=(\msk{u},\msk{\uh},\msk{\sigma},\msk{h})(r)$. $V$ evaluates the bracket without
  the mask term at $r$ from these four values, and derives the value $\gamma_A$ that $g_A(r)$ must
  take for the final check to hold.
\item[\textbf{5.}] \emph{Sumcheck B} (linearisation; $\log N$ variables, degree $2$). $V$ sends
  $\rho_B$. The parties run a sumcheck over $y\in\cube{\log N}$ for
  \[ \sum_{y}\Bigl[\,\sum_{k=1}^{3}\rho'_k(y)\,\mle{W}_k(y,r)+\frac{a_h Z(r)}{N}+\rho_B\,g_B(y)\Bigr]
     \;=\;H-\mle{z^{\mathsf{pub}}}(r)-\beta\,\mathrm{id}(r)+\alpha+\rho_B G_B ,
     \qquad \mathrm{id}(r)=\textstyle\sum_{i<\ell}2^i r_i . \]
  It ends at a point $r_y$ with final claim $\mathit{cl}_B$.
\item[\textbf{6.}] \emph{Sumcheck C} (shortness; $m$ variables, degree $4$, last round $7$). $V$ sends
  $t_z\in\F^{m}$, $k^C_1,k^C_2,k^C_3,\rho_C\in\F$. The parties run a sumcheck for
  \[ \sum_{x\in\cube{m}}\Bigl[\eq(t_z,x)\sum_{k=1}^{3}k^C_k\bigl((\msk{w}_k)^{3}-\msk{w}_k\bigr)
       +\rho_C\,g_C\Bigr](x)\;=\;\rho_C G_C . \]
  It ends at a point $r_z$. $P$ sends $\msk{w}_k(r_z)$ for $k=1,2,3$, and $V$ derives the value
  $\gamma_C$ that $g_C(r_z)$ must take.
\item[\textbf{7.}] \emph{Opening.} The protocol of Section~\ref{sec:pcs} is run for nine linear
  claims: the values $U$, $\hat U$, $s$; the final claim $\mathit{cl}_B$, which is a combination
  of $\mle{W}_k(r_y,r)$, $a_h$ and the coefficients of $g_B$; the three values $\msk{w}_k(r_z)$;
  $g_A(r)=\gamma_A$; and $g_C(r_z)=\gamma_C$.
\end{enumerate}
\end{minipage}}
\caption{The shuffle proof as a public-coin interactive protocol. In the non-interactive proof
every verifier message is a hash of the parameters, the statement and all prover messages so
far.}
\label{fig:mix}
\Description{A boxed list of the seven steps of the shuffle protocol: two commitment rounds,
three sumchecks and the opening.}
\end{figure*}

We comment on three points of the protocol.

\emph{Degrees.} The masked extensions have degree $2$ in each variable, but $Z$ vanishes as soon as
one variable is Boolean. In every round except the last, the mask terms therefore contribute
nothing to the round polynomial, and the degree is that of the unmasked summand: $3$ for
sumcheck~A and $4$ for sumcheck~C. Only the last rounds have degree $5$ and $7$.

\emph{Sumcheck B.} By the definitions of $h$ and of the masked extension,
\[
  \msk{h}(r)=\mle{z^{\mathsf{pub}}}(r)+\sum_{y}\sum_{k}\rho'_k(y)\,\mle{W}_k(y,r)
            +\beta\,\mathrm{id}(r)-\alpha+a_hZ(r),
\]
where $\mle{W}_k(y,r)$ is the extension of $W_k$ in the ballot index, evaluated at $r$ and viewed as
a table in the coefficient index $y$. The claim of sumcheck~B is this equation solved for the sum,
with the constant $a_hZ(r)$ spread over the $N$ points of the cube.

\emph{The verifier's linear algebra.} Let $A$ and $B$ be the negacyclic multiplication matrices of
$a$ and $b$. Since $\ip{\rho_u}{\coeff(ar)}=\ip{A^{\!\top}\rho_u}{\coeff(r)}$, the three blocks of
$\rho'$ are
\[
  \rho'_1=A^{\!\top}\rho_u+B^{\top}\rho_v,\qquad \rho'_2=p\,\rho_u,\qquad \rho'_3=p\,\rho_v .
\]
For the projections, the verifier needs two adjoint multiplications and $4\tau N$ multiplications of
a $\Zq$-element by an $\F$-element. This is linear in the size of the statement.

\subsection{The multiset lemma}

\begin{lemma}[Logarithmic derivative with projection]\label{lem:logder}
Let $S,T$ be multisets of $\tau<q$ elements of $\Zq^{2N}\times\Zq$. If $S\ne T$, then
\[
  \sum_{(x,y)\in S}\frac{1}{\ip{\rho}{x}+\beta y-\alpha}
  =\sum_{(x,y)\in T}\frac{1}{\ip{\rho}{x}+\beta y-\alpha}
\]
holds, with all denominators non-zero, with probability at most $(2\tau-1)/|\F|$ over uniform
$\rho\in\F^{2N}$ and $\alpha,\beta\in\F$.
\end{lemma}
\begin{proof}
Each distinct pair $(x,y)$ defines the linear form $\ip{R}{x}+By-A$ in the variables $(R,B,A)$.
Distinct pairs give distinct forms, and since every form has coefficient $-1$ at $A$, the forms are
pairwise coprime. Some pair has a net multiplicity $\mu\ne 0$ between $S$ and $T$. Since
$|\mu|\le\tau<q$, this multiplicity is non-zero in $\F$. Multiply the difference of the two sums by
the product of the $D\le 2\tau$ distinct forms. The result is a polynomial of degree at most $D-1$.
Modulo the form of the chosen pair, it equals $\mu$ times the product of the other forms, which is
non-zero. The polynomial is therefore non-zero, and the claim follows from the Schwartz--Zippel
lemma~\cite{Schwartz80,Zippel79}.
\end{proof}

The mechanised proof avoids multivariate polynomials. The numerator is homogeneous of degree
$D-1$, so on every line in a direction where it does not vanish, it is a univariate polynomial of
that degree. Averaging over the translates of the line gives the same bound. The mechanised
statement is slightly more general: the forms may have coefficients in $\F$, not only in $\Zq$. The
condition $\tau<q$ appears as ``$1,\dots,\tau$ are non-zero in $\F$''.

\subsection{Knowledge soundness}

\begin{theorem}[Knowledge soundness, interactive]\label{thm:sound}
Consider the protocol of Figure~\ref{fig:mix} with the matrices as oracles, and let $\tau<q$.
Decode every committed row, at the time its batch is committed, to the unique codeword within
distance $e$ if there is one. If the decoded $(\sigma,W)$ is not a witness for the statement, the
verifier accepts with probability at most
\[
  \frac{2\tau+9\ell+7\log N+9}{|\F|}\;+\;2\,\epsprox\;+\;\Bigl(1-\frac{e}{n}\Bigr)^{\nq}.
\]
\end{theorem}

Table~\ref{tab:budget} lists, for each verifier message, the event that lets a false claim survive
that message. The bound of the theorem is the sum of the table.

\begin{table}[t]
\caption{Error budget of Theorem~\ref{thm:sound}.}
\label{tab:budget}
\small
\begin{tabular}{@{}p{0.56\columnwidth}ll@{}}
\toprule
Event & Probability & Message\\
\midrule
log-derivative identity holds for two different multisets & $(2\tau-1)/|\F|$ & $\rho,\alpha,\beta$\\
claim of A true although an inverse or the sum is wrong & $(\ell+1)/|\F|$ & $t_c,k,\rho_A$\\
sumcheck A accepts a false claim & $(3\ell+2)/|\F|$ & rounds of A\\
$H$ wrong although the claim of B is true & $1/|\F|$ & $\rho_B$\\
sumcheck B accepts a false claim & $2\log N/|\F|$ & rounds of B\\
claim of C true although $W$ is not ternary & $(m+1)/|\F|$ & $t_z,k^C,\rho_C$\\
sumcheck C accepts a false claim & $(4m+3)/|\F|$ & rounds of C\\
a false claim disappears in the batch & $1/|\F|$ & $\kappa$\\
opening equation holds for a wrong $V$ & $1/|\F|$ & $c$\\
a far matrix yields a close combination & $\epsprox$ per matrix & $\gamma$\\
all column queries miss & $(1-e/n)^{\nq}$ & columns\\
\bottomrule
\end{tabular}
\end{table}

\begin{proof}
Assume that the verifier accepts and that none of the events of Table~\ref{tab:budget} occurs. We go
backwards through the protocol.

\emph{Opening.} Suppose a matrix had no decoding with at most $e$ bad columns. By
Lemma~\ref{lem:prox} the combination $p$ would be far from the code, so the vector sent would differ
from the committed combination on more than $e$ columns, and a query would hit one of them. Hence
both matrices decode. By Lemma~\ref{lem:cons}, every $\bar u_v$ equals $m_v+c\,u_v$ on the decoded
rows. The check of step~5 then reads $\mu'+cV'=\mu+cV$, where $\mu'$ and $V'$ are computed from the
decoded rows and are fixed before $c$ is drawn. Hence $V'=V$. Since the batching did not hide a
false claim, all nine claims hold for the decoded rows.

\emph{Sumcheck B.} The fourth claim fixes the final value of B to the true summand at $r_y$, so the
claimed sum is the true sum:
$(\msk{h}(r)-H)+\rho_B(G_B'-G_B)=0$, where $G_B'$ is the true sum of the committed $g_B$ and both
differences are fixed before $\rho_B$ is drawn. Hence $H=\msk{h}(r)$.

\emph{Sumcheck A.} Since $U,\hat U,s,H$ and $g_A(r)$ are correct, the final check of A evaluates the
true summand at $r$, so the claimed sum is true. Since $\sum_x\eq(t_c,x)=1$, the claim says that
$\sum_x\eq(t_c,x)\bigl[k_1(u h-1)+k_2(\uh\hh-1)\bigr](x)+k_3\sum_x(u-\uh)(x)$ vanishes. For random
$t_c,k_1,k_2,k_3$ this means $uh=1$ and $\uh\hh=1$ on the cube, and $\sum u=\sum\uh$. In other words,
\[
  \sum_i\frac{1}{z_i+\beta i-\alpha}=\sum_j\frac{1}{\zh_j+\beta\sigma_j-\alpha}
\]
with no vanishing denominator.

\emph{Multisets.} By Lemma~\ref{lem:logder}, the multiset of the pairs
$(\coeff(c_i)+\Mpk\coeff(w_i),\,i)$ equals the multiset of the pairs $(\coeff(\ch_j),\,\sigma_j)$. The
tags on the left are $0,\dots,\tau-1$, each exactly once, and they are distinct in $\Zq$ because
$\tau<q$. So $\sigma$ is a permutation and $\ch_j=c_{\sigma(j)}+\EncZ(w_{\sigma(j)})$.

\emph{Sumcheck C.} As for A, we get $w^3=w$ for every entry of $W$, so every entry is in
$\{-1,0,1\}$.
\end{proof}

The extractor is straight-line. It reads the committed leaves off the prover's oracle queries and
runs Reed--Solomon decoding on the first message.

\paragraph{The mechanised statement.}
Theorem~\ref{thm:sound} is machine-checked for the protocol as implemented, including the row layout
of both matrices and the verifier's computation of its public values. The bound is stated as a
count of accepting challenge tuples, for a prover that may choose every message after the
commitments adaptively; rows of the second batch may depend on the first challenge. The count splits
as in Table~\ref{tab:budget}. The sumcheck layer contributes
$2\tau+4\ell+2\log N+5m+7=2\tau+9\ell+7\log N+7$ over $|\F|$. The two opening lines for $\kappa$ and
$c$ add $2$, which gives the numerator of the theorem and confirms the paper's arithmetic. The
hypotheses are those of the theorem plus one proximity statement per matrix (Lemma~\ref{lem:prox},
cited), the Reed--Solomon parameters ($k\le n$, distinct evaluation points, $2e<d$), and three
conditions on the field: $1,\dots,\tau$, $N$ and $2$ are non-zero. They hold for $\tau,N<q$ and odd
$q$. The mask argument needs $2\ne 0$, which was not stated before. On the other hand, soundness
does not need $\rho_A,\rho_C\ne 0$. Stated without division, as
$\rho_A\,g_A(r)=\mathit{cl}_A-(\ldots)$, the claims on the masks remain sound when $\rho_A=0$. The
verifier needs $\rho_A,\rho_C\ne 0$ only in order to divide, and it rejects otherwise.

\subsection{Zero knowledge}\label{sec:zk}

We call a challenge vector \emph{good} if no $h_i$ or $\hh_j$ is zero, $Z(r)\ne 0$, $Z(r_z)\ne 0$
and $\rho_A,\rho_B,\rho_C\ne 0$. Challenges are bad with probability at most
$(2\tau+2\ell+2m+3)/|\F|$.

\begin{theorem}[Honest-verifier zero knowledge]\label{thm:zk}
There is a simulator that, given a statement in the language and good challenges, outputs a
transcript in which every value other than the hashes of unopened leaves has exactly the
distribution of a real transcript. The hashes differ from real ones with probability at most
$t/2^{128}$ for a distinguisher with $t$ oracle queries.
\end{theorem}

\paragraph{Simulator.}
(1)~Draw $U,\hat U,s,H$ and $\msk{w}_k(r_z)$ uniformly. (2)~Draw $G_A,G_B,G_C$ uniformly, and draw
every sumcheck round message uniformly among those consistent with the running claim. (3)~Derive
$\mathit{cl}_B$, $\gamma_A$, $\gamma_C$ and the batched value $V$ with the verifier's formulas.
(4)~Draw all $\bar u_v$ and $p$ uniformly and set $\mu=\sum_v\ip{\bar u_v}{L_v}-cV$. (5)~In the
opened columns, draw the entries of data rows uniformly and solve for the entries of the mask rows
from the checks of step~7. (6)~Choose salts, fill the unopened leaves with random hashes and build
the trees.

\begin{proof}
(1)~Each value is $\mle{T}(r)+a_TZ(r)$ with $a_T$ uniform and $Z\ne 0$ at the point, and each
$a_T$ is used at one point only. (2)~The non-constant coefficients of the round-$j$ polynomial of a
mask $g$ are $2^{m-j-1}c_{j,t}$, and $G$ is $2^m a_0$ plus a function of the $c_{j,t}$. The map
from the coefficients of $g$ to $(G,\text{round messages})$ is therefore a triangular affine
bijection. Because $\rho_A,\rho_B,\rho_C\neq 0$, the messages of the masked sumcheck are uniform
subject to the verifier's checks~\cite{CFS17,Libra}. (4)~The rows $m_v$ and $m_{\mathsf{prox}}$ are
uniform and each is used once, so $\bar u_v$ and $p$ are uniform, and the check of step~5 determines
$\mu$. (5)~By Lemma~\ref{lem:hide}, the opened entries of data rows are uniform; in both the real and
the simulated transcript, the checks determine the entries of the mask rows. (6)~The salts hide the
unopened leaves.
\end{proof}

\paragraph{The mechanised statement.}
Theorem~\ref{thm:zk} is machine-checked for the field values of the transcript: the masked
evaluations, the mask sums and round messages of the three sumchecks, $\mu$, the slot and
proximity vectors, and the opened entries of all rows. The statement has a permutation form. For
fixed good challenges, the list of honest transcripts over all prover coins is a permutation of the
list of simulated transcripts over all simulator coins, so the two distributions are equal. The
honest prover is instantiated with the tables, summands and row layout of the implementation. Three
points of the proof sketch above become explicit in the mechanisation.
\begin{itemize}[leftmargin=1.2em,itemsep=1pt,topsep=3pt]
\item The mask table is randomness of the sumcheck layer and, at the same time, data of a committed
  row. The map from all prover coins to the free values of the transcript is therefore not a product
  of two bijections but a triangular one: the opening's part depends on the layer's coins. A short
  composition lemma handles this.
\item The simulator of the opening needs one value from the sumcheck layer, the batched target
  value $V$ of step~(3). It recomputes $V$ from the layer transcript as the verifier does, from the
  masked evaluations and the final values of the three sumchecks. On honest transcripts this gives
  the true value.
\item The proximity vector combines all rows of a matrix, including the two $\Zq$-rows of each slot
  mask. In the opened columns, the simulator therefore solves for the slot mask rows first and for
  the proximity mask rows last.
\end{itemize}
Besides good challenges, the mechanised statement assumes that the evaluation points have order at
least $n$ and that the opened columns come from the verifier's draw. We prove that this draw gives
distinct columns, in range, and as many as there are pads. Hashes of unopened leaves are outside the
model.

\paragraph{The executable prover.}
The honest prover of the model computes its messages from functions on the cube. The executable
prover of Section~\ref{sec:mech} instead folds tables of integers and sends each round polynomial
as its values at the nodes $0,2,\dots,d$. We prove that the two agree. Folding a table at $r$
evaluates its multilinear extension at $r$ in the first free variable. Every round message of the
executable sumcheck is therefore the round polynomial of the masked summand of the original tables,
and these summands are, pointwise, the summands $g_A$, $g_B$, $g_C$ of the model. For $g_B$ this
uses the fact that the folded witness table is the extension of $W$ with its last $\log\tau$
variables fixed at $r_A$. We match the masked evaluations, the totals $G$ of the masks, $\mu$, the
slot and proximity vectors and the opened column entries in the same way, and we show that the rows
that store the masks of the opening hold the model's masks.

The result is the following. For prover coins built from the model's coins, every run of the
executable prover returns a proof whose elements of $\F$ form the model's honest transcript at the
challenges of that run. Conversely, for fixed salts, every coin record of the prover's shape with
entries in $[0,q)$ arises in this way from exactly one tuple of model coins, so uniform coins of the
prover are uniform coins of the model. Two limits of the model remain. The statement concerns the
field elements of the proof, not its hashes, salts and Merkle paths. And the challenges are fixed,
whereas in a run of the non-interactive prover they depend on the coins through the hash
(Section~\ref{sec:fs}).

%% =======================================================================================
\section{Non-interactive proofs}
\label{sec:fs}
%% =======================================================================================

We obtain the non-interactive proof with the BCS transformation~\cite{BCS16}. Oracles become Merkle
roots, and every challenge is the hash of the parameters, the statement and all earlier prover
messages. The quantity that governs its security is round-by-round soundness~\cite{Hol19}.

Round-by-round soundness requires a \emph{state function}: a predicate ``doomed'' on partial
transcripts with four properties. For a false statement the empty transcript is doomed. A complete
doomed transcript is rejected. Every prover message keeps a doomed transcript doomed or makes the
verifier reject. And every verifier message takes a doomed transcript to a non-doomed one only with
small probability, whatever the transcript so far. Our state function follows the proof of
Theorem~\ref{thm:sound}, read forwards.

As in Theorem~\ref{thm:sound}, decode every committed row when its batch is committed. Call a partial
transcript \emph{far} if one of the committed matrices has no decoding with at most $e$ bad columns.
Table~\ref{tab:rbr} assigns to each point of the protocol a \emph{claim} about the decoded rows and
the messages so far. A partial transcript is \emph{doomed} if it is far or its current claim is
false. After the challenges $c,\gamma$ of the opening, ``far'' refers to the $\gamma$-combinations
instead of the matrices. After the vectors of the opening, the transcript is doomed if a sent vector
differs on more than $e$ columns from the committed combination it stands for.

\begin{table}[t]
\caption{State function of Proposition~\ref{prop:rbr}: the claim that holds on non-doomed
transcripts after each message, and the event by which a verifier message leaves a doomed
transcript (escape). Prover messages have no escape.}
\label{tab:rbr}
\small
\begin{tabular}{@{}p{0.17\columnwidth}p{0.47\columnwidth}l@{}}
\toprule
after & claim on the decoded rows & escape\\
\midrule
commit $W,\sigma$ & $(\sigma,W)$ is a witness & --\\
$\rho,\alpha,\beta$ & $W$ ternary, and the projected multisets are equal with non-zero denominators & $\frac{2\tau-1}{|\F|}$\\
commit $u,\uh$ & also $u=1/h$, $\uh=1/\hh$ on the cube and $\sum(u-\uh)=0$ & --\\
$t_c,k,\rho_A$ & $W$ ternary, and the claim of A is true & $\frac{\ell+1}{|\F|}$\\
round of A & $W$ ternary, and the running claim of A is true & $\frac{5}{|\F|}$\\
$U,\hat U,s,H$ & $W$ ternary, and the four values and $g_A(r)=\gamma_A$ are true & --\\
$\rho_B$ & the above with $H$ replaced by the claim of B & $\frac{1}{|\F|}$\\
round of B & likewise with the running claim of B & $\frac{2}{|\F|}$\\
$t_z,k^C,\rho_C$ & the claims so far, and the claim of C & $\frac{m+1}{|\F|}$\\
round of C & likewise with the running claim of C & $\frac{7}{|\F|}$\\
$\msk{w}_k(r_z)$ & the nine linear claims of step~7 & --\\
$\kappa$ & the batched value is the value on the decoded rows ($V'=V$) & $\frac{1}{|\F|}$\\
$\mu$ & unchanged & --\\
$c,\gamma$ & $\mu'+cV'=\mu+cV$, and no $\gamma$-combination is far & $\max(\epsprox,\frac{1}{|\F|})$\\
vectors & every vector agrees with its committed combination on all but $e$ columns & --\\
columns & every queried column agrees & $(1-\frac{e}{n})^{\nq}$\\
\bottomrule
\end{tabular}
\end{table}

\begin{proposition}[Round-by-round soundness]\label{prop:rbr}
With the state function of Table~\ref{tab:rbr}, the interactive protocol has round-by-round
knowledge error
\[
  \epsrbr=\max\Bigl\{\frac{\max(2\tau-1,\;m+1,\;7)}{|\F|},\;\epsprox,\;\Bigl(1-\frac{e}{n}\Bigr)^{\nq}\Bigr\}
  \;\le\;2^{-128}
\]
for our parameters.
\end{proposition}

\begin{proof}
\emph{Initial and final states.} If the decoded $(\sigma,W)$ is not a witness, the empty
transcript after the first commitment is doomed. A complete doomed transcript is rejected: either a
queried column differs from what the vectors imply, or the opening equation of step~5 of the opening
fails.

\emph{Verifier messages.} Each row of Table~\ref{tab:rbr} with an escape is one event of
Table~\ref{tab:budget}, now conditioned on a fixed doomed prefix. For $\rho,\alpha,\beta$ it is
Lemma~\ref{lem:logder}; ternarity does not depend on the challenge. For $t_c,k,\rho_A$, for $\rho_B$
and for $t_z,k^C,\rho_C$, the claim after the message is a random linear combination, or the value of
a polynomial of degree at most $\ell+1$, $1$ and $m+1$ respectively, of quantities fixed before the
message, one of which is non-zero. In every sumcheck round, a false running claim becomes true for
at most $d$ challenges, whatever round polynomial the prover sends, with $d\le 5$, $2$, $7$. For
$\kappa$, a non-zero vector of claim errors gives a zero combination with a uniform $\kappa$ with
probability $1/|\F|$. For $c,\gamma$, the two escapes belong to disjoint cases. If no matrix is far,
the equation $\mu'-\mu=c(V-V')$, with $V\ne V'$ fixed before $c$, has one solution. If a matrix is
far, its $\gamma$-combination is close with probability at most $\epsprox$ (Lemma~\ref{lem:prox}).
For the columns, a vector that differs from its committed combination on more than $e$ of the $n$
columns is caught unless all $\nq$ queries miss.

\emph{Prover messages.} Committing more rows cannot undo ``far''. The second commitment keeps a
false multiset equality false: if $u=1/h$ and $\uh=1/\hh$, then $\sum(u-\uh)$ is the difference of the
two sides of Lemma~\ref{lem:logder}. The values $U,\hat U,s,H$ and $\msk{w}_k(r_z)$ either fail the
final check of their sumcheck, and the verifier rejects, or one of the claims they introduce is
false, since the final check equates the running claim with the summand evaluated from them. The
value $\mu$ does not enter the claim of its row. The vectors of the opening either agree with the
committed combinations on all but $e$ columns, in which case the opening equation is the equation
on the decoded rows by Lemma~\ref{lem:cons}, or they do not, in which case the next state is doomed
by definition.

Since $\ell+1\le m+1$ and $5\le 7$, the largest escape is the bound of the proposition.
\end{proof}

This proof is on paper. Its transitions are the counted events of the mechanisation. The sumcheck
transition is mechanised in exactly this form (\texttt{SumcheckSound.v: sc\_round\_rbr}: from a
doomed state, at most $d$ challenges lead to a state that is not doomed, for any round polynomial;
\texttt{sc\_doomed\_final}), and the other escapes are the per-message counts used in the proof of
Theorem~\ref{thm:sound}. The state function itself, the composition over the messages, and the BCS
transformation with its Merkle trees are not mechanised.

\begin{theorem}[Non-interactive security]\label{thm:nizk}
Assume Proposition~\ref{prop:rbr}. In the random-oracle model, the non-interactive shuffle proof
has knowledge error at most $t\,\epsrbr+O(t^2/2^{256})$ against provers with $t$ classical
queries~\cite{BCS16,Hol19}, and $O(t^2\,\epsrbr+t^3/2^{256})$ against provers with $t$ quantum
queries~\cite[Thm.~3]{CMS19}. It is statistically zero-knowledge.
\end{theorem}

The quantum statement assumes that the theorem of~\cite{CMS19} applies to our concrete hashing
layout, which differs in details from the textbook compiler: salted column leaves, a hash-chain
transcript, and rejection sampling of field elements. With $\epsrbr\approx 2^{-128}$, the quantum
bound is meaningful for up to about $2^{64}$ queries. A deployment that targets a higher quantum
query bound would increase $\nq$.

%% =======================================================================================
\section{The decryption proof}
\label{sec:dec}
%% =======================================================================================

A decryption server proves knowledge of a witness for $\Rdec$. Let $B=\mathit{base}^{\,k}$ with
$B<q$. The prover writes every coefficient of the noise as
\[
  E=\sum_{d<k}\mathit{base}^{\,d}\,\delta_d-B/2,\qquad \delta_d\in\{0,\dots,\mathit{base}-1\},
\]
and the protocol runs as follows.
\begin{enumerate}[leftmargin=1.5em,itemsep=1pt,topsep=3pt]
\item The prover commits to the digit tables $\delta_0,\dots,\delta_{k-1}$, each of length $\tau N$,
  in the big matrix, and to $s$, $e$ and the mask table in the small matrix.
\item The verifier sends $\theta\in\F^{\ell}$ and $\theta_1,\theta_2\in\F^{\log N}$. This compresses
  relations (D) and (K) into two linear claims with public value zero:
  \begin{align*}
    \sum_{i,x}\eq(\theta,i)\,\eq(\theta_1,x)\,\bigl[t_i-su_i-pE_i\bigr]_x&=0,\\
    \sum_{x}\eq(\theta_2,x)\,\bigl[b-as-pe\bigr]_x&=0 ,
  \end{align*}
  where $[\cdot]_x$ denotes coefficient $x$. The product $su_i$ is linear in $s$, and, as in the
  shuffle proof, the verifier moves it to the public side with the adjoint of multiplication by
  $u_i$.
\item A masked sumcheck over $m$ variables, of degree $\mathit{base}+1$ and of degree
  $2\,\mathit{base}+1$ in the last round, proves $\prod_{v<\mathit{base}}(\delta_d-v)=0$ for all
  digits.
\item A masked sumcheck over $\log N$ variables, of degree $4$ and of degree $7$ in the last round,
  proves $s^3=s$ and $e^3=e$.
\item One opening settles $k+6$ linear claims.
\end{enumerate}

\begin{theorem}\label{thm:dec}
The decryption proof is honest-verifier zero knowledge in the sense of Theorem~\ref{thm:zk}. As an
interactive oracle proof, it has knowledge error at most
\[
  \frac{(\mathit{base}+3)\,m+6\log N+\mathit{base}+7}{|\F|}+2\,\epsprox+\Bigl(1-\frac{e}{n}\Bigr)^{\nq}.
\]
\end{theorem}

The proof follows those of Theorems~\ref{thm:sound} and~\ref{thm:zk}. The range proof is
\emph{exact}: digits in $\{0,\dots,\mathit{base}-1\}$ and $B<q$ give $-B/2\le E<B/2$ over the
integers, so there is no gap between the bound that honest provers satisfy and the bound that is
proven. Relation (K) with ternary $s,e$ is also what a server proves about its key share at key
generation.

The soundness half of Theorem~\ref{thm:dec} is machine-checked for the protocol as implemented, with
the same structure and hypotheses as for the shuffle proof. The sumcheck layer contributes
$(\mathit{base}+3)\,m+6\log N+\mathit{base}+5$ over $|\F|$, and the two opening lines bring this to the
numerator of the theorem. The adjoint of negacyclic multiplication connects the two compressed
claims to the committed rows: the public side minus the claim's value on the decoded rows is the
multilinear extension of the defect table of (D) or (K), respectively. For the range proof, the
mechanisation shows that every digit lies in $\{0,\dots,\mathit{base}-1\}$ and that such digits
give an integer below $B$. Every noise coefficient is therefore $v-B/2$ for a natural number $v<B$,
that is, $-B/2\le E<B/2$.

Zero knowledge of the decryption proof is machine-checked in the same form as for the shuffle proof
(Theorem~\ref{thm:zk}): for fixed good challenges, the honest transcripts over all coins are a
permutation of the simulated ones. The sumcheck layer, with its two masked zero-test sumchecks, is
proven for any field. The opening holds for any committed data, because its check equation is the
batching identity together with the fact that pads lie outside every slot vector. The witness enters
only where the simulator of the opening recomputes the batched target value from the layer
transcript, which uses the two compressed claims and the row layout.

The executable decryption prover refines this model in the same way as the shuffle prover
(Section~\ref{sec:zk}). For prover coins built from the model's coins, every run returns the model's
honest transcript at its challenges, and coin records of the prover's shape correspond one to one to
model coins. Here the mask table of the layer is the third data row of the small matrix, and the
random rows of both matrices hold the real and imaginary parts of the slot and proximity masks.

%% =======================================================================================
\section{Mechanisation}
\label{sec:mech}
%% =======================================================================================

The development is written in the Rocq prover~\cite{Rocq}, version~9.2. It has about $26{,}500$
lines of its own, plus about $11{,}000$ lines of general libraries for algebra, probability and
utilities. We check that $335$ main theorems are closed under the global context, that is, that they
use no axioms. The primality certificate of $q$ is checked with the Pocklington checker of
CoqPrime~\cite{GTW06}, which needs no axioms either.

\paragraph{Structure.}
The generic part (about $13{,}100$ lines) works over an arbitrary field with decidable equality. It
contains the multilinear extension, sumcheck soundness and completeness, the masked sumcheck,
Reed--Solomon codes and the opening, the row layouts of both proofs, the two soundness theorems, and
the zero-knowledge statements for transcripts. Probabilities are counts. Given a duplicate-free list
of all field elements, an event has probability at most $k/|\F|$ when at most $k\,|\F|^{n-1}$
challenge tuples satisfy it, and the composition of several messages is a lemma about such counts.
The instance part (about $13{,}200$ lines) fixes $\F=\Zq[t]/(t^2-3)$ for the $q$ of
Section~\ref{sec:prelim}, proves that it is a field, and instantiates the generic theorems. It also
contains the executable code.

\paragraph{From the protocol to the implementation.}
Two decisions make the soundness theorems apply to the code rather than to an idealised protocol.
First, we model the opening as the implementation runs it. Each claim is a list of terms (slot, row,
coefficient, public vector), tensor and functional slots are compiled into the vectors and rows that
the verifier uses, and a mask row over $\F$ is a pair of $\Zq$-rows with weights $1$ and $t$.
Second, the row layout is explicit: which rows hold $W$, $\sigma$, $u$, $\uh$ and the mask table, and
at which offsets the mask table stores each mask. For this layout we prove that the values of the
claims on the decoded rows equal the ideal values used by the sumcheck argument.

\paragraph{Executable code.}
The verifiers and provers are written in Rocq over pairs of integers, and we prove that they compute
the field operations of $\F$. The extraction to OCaml~\cite{Letouzey02} maps integers to arbitrary
precision. Each verifier has a \emph{front end}, which parses the bytes, replays the Fiat--Shamir
transcript and checks the Merkle paths, and a \emph{back end}, which does the algebra. For the back
end we prove the following: if it accepts, then, for the challenges of the transcript, the claimed
values, the claims and the opening checks are exactly the events whose number the soundness theorems
bound.

For the shuffle prover we prove completeness: for a valid witness and any random coins, if the prover
returns a proof, the verifier accepts it. The prover returns no proof only when it would divide by
zero. For zero knowledge we prove that the prover refines the model of Theorem~\ref{thm:zk}
(Section~\ref{sec:zk}). The decryption prover is written in the same way and reuses the claims and
slots of the decryption verifier. We prove it complete as well, and we prove that it refines the
zero-knowledge model of the decryption proof in the same sense as the shuffle prover. Its digit check
is the product $\prod_{v<\mathit{base}}(D-v)$ expanded into its coefficients. Besides relations (D)
and (K) and ternary $s,e$, completeness needs $(E+\mathit{off}) \bmod q < \mathit{base}^{k}$, so that
the digits recombine to the noise. The digit check itself passes for any noise, because every digit
is reduced modulo $\mathit{base}$.

SHA-256 is HACL*'s verified implementation~\cite{HACL17}, called through a thin OCaml binding. We also
keep a Rocq definition of SHA-256 and compare it with HACL*'s in tests.

\paragraph{Trusted base.}
The following are not proven and are trusted when the extracted code runs: the parsers, the
transcript and the Merkle checks, which are only tested; the OCaml drivers that read files and
compute code parameters; the extraction mechanism, the OCaml compiler and the integer library; the
replacement of ten list functions and four arithmetic functions on natural numbers from the standard
library by their OCaml counterparts; and the agreement of HACL*'s SHA-256 with the Rocq definition,
which is tested. The soundness theorems themselves are about the interactive protocols, and
Section~\ref{sec:fs} is needed to transfer them to the hashed proofs.

\paragraph{Engineering notes.}
Two pitfalls cost us time, and we record them for similar developments. A proof about code that
draws Fiat--Shamir challenges should keep the drawing operation abstract. With the concrete
challenger, a case analysis on a draw can make the kernel evaluate SHA-256 on symbolic data while it
checks the proof. Similarly, a statement should not hide the enumeration of all elements of $\F$
behind a defined constant. A conversion that unfolds the constant starts computing a list of $q^2$
elements. One such proof took eight minutes to check before we changed it, and about a second after.

\subsection{What the formalisation found}
\label{sec:findings}

\begin{itemize}[leftmargin=1.2em,itemsep=2pt,topsep=3pt]
\item The error bounds of Theorems~\ref{thm:sound} and~\ref{thm:dec} are confirmed term by term.
\item The masked sumcheck needs $2\ne 0$ in $\F$, and the sumcheck that linearises $h$ needs
  $N\ne 0$. Both hold in odd characteristic $q>N$, but neither condition had been stated.
\item Soundness does not need $\rho_A,\rho_C\ne 0$ (Section~\ref{sec:shuffle}).
\item The opening has two kinds of slot, and the earlier description covered only one
  (Section~\ref{sec:pcs}). For two matrices, the query term is a maximum, not a sum.
\item The proof of Appendix~\ref{app:prox} holds with the constant $(e+2)/|\F|$. We restructured its
  second step as an induction on the rows (decoding is affine), followed by the covering argument.
\item The extracted witness lives in $\F$. Soundness does not need the fact that the committed rows
  are over $\Zq$.
\item The zero-knowledge argument needs the composition and ordering details listed after
  Theorem~\ref{thm:zk}. The simulator of the implementation already followed them.
\end{itemize}

%% =======================================================================================
\section{Evaluation}
\label{sec:eval}
%% =======================================================================================

All measurements in this section run code extracted from the development of
Section~\ref{sec:mech}: the shuffle prover and the shuffle verifier, each with its unverified OCaml
driver, which reads the files and computes the code parameters. The extracted code is written to be
proven correct, not to be fast: it computes with lists of arbitrary-precision integers.

\paragraph{Setup.}
We use $N=4096$, $q\approx 2^{78}$, $\lambda=128$, proximity regime $\Delta=3$ and rate $1/4$, on one
core of an Apple M5~Max with OCaml~5.4, Zarith, and SHA-256 from HACL*. The prover draws its coins
from the operating system. A separate script generates statements and witnesses, and we do not time
this step.

\begin{table}[t]
\caption{Extracted shuffle prover and verifier, one core, $N=4096$, $\lambda=128$, $\Delta=3$,
rate $1/4$. Memory is the peak resident size.}
\label{tab:xbench}
\small
\begin{tabular}{@{}rrrrrr@{}}
\toprule
ballots & prove & verify & proof & prover & verifier\\
 & (s/ballot) & (s) & (MB) & memory (GB) & memory (GB)\\
\midrule
16  & 0.43 & 1.0 & 0.90 & 0.46 & 0.20\\
256 & \xbProveA & \xbVerA & \xbSizeA & 5.9 & 2.1\\
1024 & 0.29 & 17.6 & 3.78 & 19.8 & 8.3\\
\bottomrule
\end{tabular}
\end{table}

\paragraph{Results.}
Table~\ref{tab:xbench} gives the results. The extracted verifier accepts the proof in every run. The
proof size follows from the parameters: it grows with the square root of the number of ballots
(Section~\ref{sec:overview}), while the statement grows linearly. The prover's time per ballot falls
as the number of ballots grows, because part of its work is per proof. Its memory is dominated by the
witness and the encoded matrices, which are held as lists of boxed integers.

Three changes made the extracted code usable at this size. First, the verifier projects the public
key with the adjoint of negacyclic multiplication (Section~\ref{sec:shuffle}), and the prover needs
the same projection. Computed directly, it takes $N^2$ products. We replaced it by three NTTs and a
pointwise product, and proved that the result is the same: at the roots $z$ of $X^N+1$, the adjoint
of $a$ applied to $\rho$ takes the value $a(z^{-1})\rho(z)$, and an inverse NTT recovers the
coefficients. The code checks the required facts about the chosen root of unity at run time and
falls back to the direct computation if they fail, so the equality holds without hypotheses.
Extracted code does not share values between definitions, and the prover computed this projection
$21$ times per proof. At $16$ ballots the change took the prover from $209$\,s to $7$\,s, and the
verifier from $10.6$\,s to $1.0$\,s. Second, at extraction, the list functions of the standard
library (\texttt{map}, \texttt{app}, \texttt{seq} and a few others) are replaced by OCaml's
tail-recursive versions. The extracted definitions recurse once per element, and OCaml~5 scans the
whole stack at every minor collection, which made the verifier quadratic in practice. Third, the
opened columns of the slot and proximity vectors are read from one NTT encoding per vector instead of
one Horner evaluation per column. Again we prove that the check computes the same values.

\paragraph{Tests.}
On four parameter sets in each proximity regime ($N$ from $8$ to $64$, $2$ to $64$ ballots), the
extracted verifier accepts the proofs of the extracted prover. It rejects ten bit-flipped proofs, a
truncated proof, and the statement in which two outputs exchange one CRT slot
(Section~\ref{sec:overview}). We compare a Rocq definition of SHA-256 with HACL*'s on test inputs.

\paragraph{Decryption proof.}
For $16$ ballots, $28$ base-$4$ digits (a noise range of $2^{56}$) and the parameters above, the
extracted decryption prover takes $62.5$\,s ($3.9$\,s per ballot, $2.4$\,GB) and the extracted
decryption verifier takes $2.5$\,s. The proof has $1.93$\,MB. On four parameter sets in each
proximity regime (digit bases $2$, $4$, $8$ and $16$), the extracted verifier accepts the extracted
prover's proofs and rejects six bit-flipped proofs, a truncated proof and a changed statement.
Running the prover at full size exposed a cost in the verifier's claims. The digit weights
$\mathit{base}^d$ were computed as natural numbers, which extraction implements by repeated addition:
$2^{54}$ steps for $d=27$. They are now computed as integers, with the same proof.

%% =======================================================================================
\section{Related work}
\label{sec:related}
%% =======================================================================================

\paragraph{Lattice shuffles.}
Costa, Mart\'inez and Morillo~\cite{CMM19} and Strand~\cite{Strand18} gave early proofs of shuffle
for lattice ciphertexts. Aranha et al.~\cite{ABGSctrsa21} and Farzaliyev, Willemson and
Kaasik~\cite{FWK21} gave practical constructions for voting, and~\cite{ABGS23} gave the mix-net with
distributed decryption that we build on. Herranz, Mart\'inez and S\'anchez~\cite{HMS21} obtain
sublinear lattice-based shuffle proofs by expressing the shuffle as an arithmetic circuit over $\Zq$;
their work has no implementation. Bootle et al.~\cite{BLM25} found the soundness flaw discussed above
and repaired the lattice-based argument with a product argument whose tags have invertible pairwise
differences. Our tagged logarithmic derivative plays the same role over a field. Abdolmaleki et
al.~\cite{AbdShuffle} give a lattice-based shuffle argument from subset checking, and
Orthus~\cite{Orthus} batch-verifies lattice relations with sublinear proofs. Our work differs from all
of these in the proof system: we use no lattice commitments, and we obtain statistical soundness in
the random-oracle model. We do not know of an earlier hash-based proof of shuffle for lattice
ciphertexts, but our literature search was limited.

\paragraph{Verified mix-nets and cryptographic code.}
Haines, Gor\'e and Sharma~\cite{HGS21} verified in Coq a verifier for the discrete-logarithm shuffle
of Terelius and Wikstr\"om and extracted it to check real election data. Our mechanisation follows a
similar route for a hash-based proof over lattice ciphertexts, and it also covers zero knowledge and
the completeness of the provers. HACL*~\cite{HACL17} provides the verified SHA-256 that our extracted
code calls.

\paragraph{Proof systems.}
The commitment is that of Ligero~\cite{AHIV17}, used as a polynomial commitment as in
Brakedown~\cite{Brakedown}, and the proximity analyses are from~\cite{AHIV17,DP23,BCIKS20}. The masked
sumcheck is from~\cite{CFS17,Libra}, and the logarithmic-derivative argument is from
Hab\"ock~\cite{Hab22}. To our knowledge, the combination of a random projection of ring elements with
a tagged logarithmic derivative is new, and so is the resulting separation of the linear part of the
relation from the permutation and shortness conditions. Succinct lattice-based
commitments~\cite{LaBRADOR,NS24} or a folding opening~\cite{BaseFold} could replace our opening and
make the proof smaller, at the price of computational soundness or a more involved analysis.

\paragraph{Compressing ring elements to a field.}
Compressing ring elements to a field with a random functional has precedents. Fiore, Gennaro and
Pastro~\cite{FGP14} hash BV ciphertexts by evaluating them, without reduction modulo the cyclotomic
polynomial, at a random point. This hash is a ring homomorphism with collision probability $(c+N)/q$.
Fiore, Nitulescu and Pointcheval~\cite{FNP20} handle the reduction modulo $R(X)$ with a committed
quotient and evaluation at a random point. These hashes preserve multiplication, which our relation
does not need: it multiplies only by the public key, so it is $\Zq$-linear in the witness. We
therefore use an unstructured random linear form over $\mathrm{GF}(q^2)$, whose collision probability
$1/|\F|$ does not depend on $N$. The adjoint step $\rho'=\rho\Mpk$ uses the identity behind lattice
proofs of inner products: the inner product of coefficient vectors is the constant coefficient of
$\sigma_{-1}(a)\,b$, where $\sigma_{-1}$ maps $X$ to $X^{-1}$, as in Lyubashevsky, Nguyen and
Plan\c{c}on~\cite{LNP22}.

%% =======================================================================================
\section{Limitations and future work}
\label{sec:limits}
%% =======================================================================================

\paragraph{Proofs.}
The soundness and zero-knowledge theorems of both proofs and Appendix~\ref{app:prox} are
machine-checked. The following are not. The line lemma and the proximity gap
of~\cite{DP23,BCIKS20} are cited and enter the mechanisation as hypotheses; we checked the cited
statements against the papers, but not page by page. Proposition~\ref{prop:rbr} is proven on paper:
its transitions are mechanised, but its state function and the composition are not. The BCS
transformation with its Merkle trees is outside the mechanised model, and Theorem~\ref{thm:nizk}
relies on our hashing layout matching the compiler analysed in~\cite{BCS16,CMS19}. The
zero-knowledge theorems are mechanised for provers over $\F$ built from the tables and row layout of
the executable code. For both proofs, we prove that the extracted prover sends exactly these
messages, at fixed challenges and apart from hashes, salts and Merkle paths. Finally, the theorems
count the challenges for which events occur. The step from these counts to the success probability
of an adversary in the random-oracle model is the standard one and is not mechanised.

\paragraph{Scope.}
We prove statements about one mix server and one decryption server. We do not give a statement about
the whole mix-net, including the noise budget across several mix servers and the privacy that follows
from BGV. We do not analyse key generation. Relation (K) lets a server prove that its key share is
well formed, but a distributed key generation that is secure against rushing adversaries is future
work. The number of ballots must be a power of two. The natural fix is to pad with trivial
ciphertexts at public positions, which needs a short lemma.

\paragraph{Implementation.}
The extracted code computes with lists of arbitrary-precision integers and is slow
(Section~\ref{sec:eval}). A version over machine words with arrays, proven equal to the present one,
is future work. The proof format hashes leaves, nodes and the transcript without domain separation,
and the prover is not constant-time. The drivers that read files and compute the code parameters are
not verified. These issues must be fixed before any use beyond experiments.

\paragraph{Smaller proofs.}
We have estimated, but not built, two restructurings. Folding the linearisation into the shortness
sumcheck would open the big matrix at one point instead of two, for an estimated reduction of $14\%$
to $20\%$. Using the slot vector as the proximity test would then give an estimated $35\%$ to $41\%$
in regime $\Delta=3$, and requires the tensor proximity result of~\cite{DP23}.

%% =======================================================================================
\section{Conclusion}
%% =======================================================================================

Apart from a permutation and a shortness condition, the mix-net relation for BGV is linear over
$\Zq$. We project each ciphertext to one field element and prove the relation over a field. This
removes the algebraic structure that broke the ring-based shuffle argument and lets us use standard
hash-based tools; neither proof rests on a lattice assumption. The soundness and zero knowledge of
both proofs, and the correctness of the executable verifiers and provers, are machine-checked. Both
extracted provers are complete and send the messages of the zero-knowledge model. Only a cited
proximity lemma and the Fiat--Shamir step are left to pen and paper.

\bibliographystyle{ACM-Reference-Format}
\bibliography{refs}

\appendix

%% =======================================================================================
\section{From lines to affine spaces}
\label{app:prox}
%% =======================================================================================

We prove Lemma~\ref{lem:prox} for $\Delta=3$ from the line lemma. Let $\mathcal{C}$ be a linear
code of length $n$ and distance $d$ over $\F$, and let $e\le\lfloor(d-1)/3\rfloor$ with $e+1<|\F|$.
A word is \emph{close} if it is within distance $e$ of $\mathcal{C}$.

\begin{lemma}[Line lemma {\cite[Thm.~2.1]{DP23}}]\label{lem:line}
Let $x,y\in\F^n$. If more than $e+1$ points $\gamma\in\F$ make $x+\gamma y$ close, then there is a
set of at most $e$ coordinates outside which both $x$ and $y$ agree with codewords.
\end{lemma}

Let $U=m_{\mathsf{prox}}+\mathrm{span}(\mathit{row}_1,\dots,\mathit{row}_R)$, parametrised by
$\gamma\in\F^{R}$, and suppose that more than an $(e+2)/|\F|$ fraction of $U$ is close. If $U$ is a
line, Lemma~\ref{lem:line} applies directly, so assume $\dim U\ge 2$.

\emph{Every point of $U$ is close.} Fix $u\in U$. The lines through $u$ partition
$U\setminus\{u\}$, and each of them contains $|\F|-1$ points other than $u$. If every such line had
at most $e+1$ close points other than $u$, the fraction of close points in $U$ would be at most
$(e+1)/(|\F|-1)\cdot(1-1/|U|)+1/|U|\le(e+2)/|\F|$. So some line through $u$ has at least $e+2$ close
points other than $u$. Lemma~\ref{lem:line}, applied to a base point and the direction of that line,
shows that every point of the line is close, in particular $u$.

\emph{A common error set.} Any two points of $U$ lie on a line whose $|\F|>e+1$ points are all close.
By Lemma~\ref{lem:line} they agree with codewords outside a common set of at most $e$ coordinates.
Since $2e<d$, decoding within distance $e$ is unique, so the decoding map is affine on $U$. The error
vectors of the points of $U$ therefore form an affine space in which every vector has weight at most
$e$. Suppose its joint support had more than $e$ coordinates. Then some combination of at most $e+1$
of its elements, with coefficients that avoid at most $e+1$ bad values, would have weight above $e$;
such coefficients exist because $e+1<|\F|$. Hence the joint support has at most $e$ coordinates.
Outside it, every element of $U$ and every direction of $U$, in particular every row, agrees with a
codeword. This contradicts the hypothesis of Lemma~\ref{lem:prox}.

\end{document}
